\subsection{群上的中心滤过}

定义 2. 设 G 是一个群。G 上的实滤过 \((\mathrm{G}_{\alpha})\) 称为中心的，如果 \(\mathrm{G} = \bigcup_{\alpha >0}\mathrm{G}_{\alpha}\) ，并且换位子 \((x,y) = x^{-1}y^{-1}xy\) ——它由一个元素 \(x\) （属于 \(\mathrm{G}_{\alpha}\) ）与一个元素 \(y\) （属于 \(\mathrm{G}_{\beta}\) ）构成——属于 \(\mathrm{G}_{\alpha +\beta}\) 。

用阶函数 v 来说，上述定义转换为关系式

\[
v (x) > 0, \quad v ((x, y)) \geqslant v (x) + v (y) \quad \text { for   all } x, y \text { in } G.\tag{10}
\]

我们推出 \(v((x,y)) > v(x)\) （若 \(v(x) \neq +\infty\) ）；若记 \(x^y = y^{-1}xy\) （参见《代数》， 第 I 章， § 6， no. 2），则 \(x^y = x.(x,y)\) ，从而

\[
v (x ^ {y}) = v (x).\tag{11}
\]

这个关系式表达了如下事实：G 的诸子群 \(G_{\alpha}\) 中的每一个都是正规的。诸 \(G_{\alpha}\) 构成与 G 上群结构相容的一个拓扑 (《一般拓扑学》， 第 III 章， § 1， no. 2， 例) 之下 e 的一个邻域基本系，该拓扑称为由滤过 ( \(G_{\alpha}\) ) 所定义。

在本 no. 的其余部分，G 表示一个带中心滤过 \((G_{\alpha})\) 的群。对一切 \(\alpha \in R\) ，我们用下式定义 G 的子群 \(G_{\alpha}^{+}\) ：

\[
\mathrm{G} _ {\alpha} ^ {+} = \bigcup_ {\beta > \alpha} \mathrm{G} _ {\beta}.\tag{12}
\]

特别地， \(G_{\alpha}^{+} = G_{\alpha} = G\) （对 \(\alpha \leqslant 0\) ）。回顾若 A 与 B 是 G 的两个子群，则 (A, B) 表示由换位子 \((a, b)\) （满足 \(a \in A\) 且 \(b \in B\) ）生成的 G 的子群。用这个记号，我们有公式

\begin{enumerate}
\def\labelenumi{(\arabic{enumi})}
\setcounter{enumi}{12}
\tightlist
\item
\end{enumerate}

\[
\left(\mathrm{G} _ {\alpha}, \mathrm{G} _ {\beta}\right) \subset \mathrm{G} _ {\alpha + \beta}\tag{13}
\]

\[
\left(\mathrm{G} _ {\alpha} ^ {+}, \mathrm{G} _ {\beta}\right) \subset \mathrm{G} _ {\alpha + \beta} ^ {+}\tag{13'}
\]

\[
(G, G _ {\alpha}) \subset G _ {\alpha} ^ {+}.\tag{14}
\]

由 (14)， \(\mathrm{G}_{\alpha}^{+}\) 是 \(\mathrm{G}_{\alpha}\) 的一个正规子群（对一切 \(\alpha \in \mathbf{R}\) ），并且商群 \(\mathrm{gr}_{\alpha}(\mathrm{G}) = \mathrm{G}_{\alpha} / \mathrm{G}_{\alpha}^{+}\) 是交换的。我们记 \(\mathrm{gr}(\mathrm{G}) = \bigoplus_{\alpha \in \mathbb{R}}\mathrm{gr}_{\alpha}(\mathrm{G})\) ，并给这个群赋予 \(\mathbf{R}\) 型分次，其中 \(\mathrm{gr}_{\alpha}(\mathrm{G})\) 由 \(\alpha\) 次的元素组成。于是 \(\mathrm{gr}_{\alpha}(\mathrm{G}) = \{0\}\) （对 \(\alpha \leqslant 0\) ）。

命题 1. (i) 设 \(\alpha, \beta\) 属于 \(\mathbf{R}\) 。存在一个双加性映射

\[
\phi_ {\alpha \beta}: \operatorname{gr} _ {\alpha} (\mathrm{G}) \times \operatorname{gr} _ {\beta} (\mathrm{G}) \rightarrow \operatorname{gr} _ {\alpha + \beta} (\mathrm{G})
\]

它把 \((x\mathrm{G}_{\alpha}^{+},y\mathrm{G}_{\beta}^{+})\) 映到 \((x,y)\mathrm{G}_{\alpha +\beta}^{+}\) 上

\begin{enumerate}
\def\labelenumi{(\roman{enumi})}
\setcounter{enumi}{1}
\item
  设 \(\phi\) 是 \(\operatorname{gr}(\mathbf{G}) \times \operatorname{gr}(\mathbf{G})\) 到 \(\operatorname{gr}(\mathbf{G})\) 中的一个双加性映射，它在 \(\operatorname{gr}_{\alpha}(\mathbf{G}) \times \operatorname{gr}_{\beta}(\mathbf{G})\) 上的限制是 \(\phi_{\alpha \beta}\) （对每个有序对 \((\alpha, \beta)\) ）。映射 \(\phi\) 给 \(\operatorname{gr}(\mathbf{G})\) 一个李 Z-代数结构。
\item
  回顾恒等式
\end{enumerate}

\[
(x x ^ {\prime}, y) = (x, y) ^ {x ^ {\prime}} \left(x ^ {\prime}, y\right)\tag{15}
\]

其中 \(x, x', y\) 属于 \(\mathbf{G}\) (《代数》， 第 I 章， § 6， no. 2， 公式 (4 bis))。

对 \(x \in G_{\alpha}\) 与 \(y \in G_{\beta}\) ，模 \(G_{\alpha + \beta}^{+}\) 的类（元素 \((x, y)\) 取自 \(G_{\alpha + \beta}\) ）将记作 \(f(x, y)\) 。对 \(a\) （属于 \(G_{\alpha + \beta}\) ）与 \(x'\) （属于 \(G, a^{-1}.a^{x'} = (a, x') \in G_{\alpha + \beta}^{+}\) ）；特别地， \(f(x, y)\) 等于模 \(G_{\alpha + \beta}^{+}\) 的类（元素 \((x, y)^{x'}\) 的类）。因此公式 (15) 蕴涵

\[
f (x x ^ {\prime}, y) = f (x, y) f (x ^ {\prime}, y).\tag{16}
\]

而 \((y, x) = (x, y)^{-1}\) ，从而

\[
f (y, x) = f (x, y) ^ {- 1}.\tag{17}
\]

由 (16) 与 (17) 我们推出

\[
f (x, y y ^ {\prime}) = f (x, y) f (x, y ^ {\prime}).\tag{18}
\]

我们必须证明映射 \(f\colon \mathrm{G}_{\alpha}\times \mathrm{G}_{\beta}\to \mathrm{gr}_{\alpha +\beta}(\mathrm{G})\) 在取商后定义一个映射 \(\phi_{\alpha \beta}\colon \mathrm{gr}_{\alpha}(\mathrm{G})\times \mathrm{gr}_{\beta}(\mathrm{G})\to \mathrm{gr}_{\alpha +\beta}(\mathrm{G})\) 。由 (16) 与 (18)，只需证明 \(f(x,y) = 0\) （当 \(x\in \mathrm{G}_{\alpha}^{+}\) 或 \(y\in \mathrm{G}_{\beta}^{+}\) 时），而这由 (13') 得出。

\begin{enumerate}
\def\labelenumi{(\roman{enumi})}
\setcounter{enumi}{1}
\tightlist
\item
  由于 \((x, x) = e\) ，由 (17) 推出 \(\phi\) 是一个交错的双 \(\mathbf{Z}\) -线性映射。因此剩下要证明的是，对 \(u \in \mathrm{gr}_{\alpha}(\mathrm{G})\) 、 \(v \in \mathrm{gr}_{\beta}(\mathrm{G})\) 与 \(w \in \mathrm{gr}_{\gamma}(\mathrm{G})\) ，
\end{enumerate}

\[
\phi (u, \phi (v, w)) + \phi (v, \phi (w, u)) + \phi (w, \phi (u, v)) = 0.\tag{19}
\]

设 \(x \in G_{\alpha}, y \in G_{\beta}\) 与 \(z \in G_{\gamma}\) 是分别表示 \(u, v\) 与 \(w\) 的元素。我们知道 \(x^{y}\) 与 \(x\) 是 \(G_{\alpha}\) 的两个模 \(G_{\alpha}^{+}\) 同余的元素，因而 \(x^{y}\) 是 \(u\) 在 \(G_{\alpha}\) 中的一个代表；由于 \((y, z)\) 是 \(\phi(v, w)\) 在 \(G_{\beta + \gamma}\) 中的一个代表，我们看出 \((x^{y}, (y, z))\) 是 \(\phi(u, \phi(v, w))\) 在 \(G_{\alpha + \beta + \gamma}\) 中的一个代表。通过轮换，我们看出 \((y^{z}, (z, x))\) 与 \((z^{x}, (x, y))\) 分别表示 \(\phi(v, \phi(w, u))\) 与 \(\phi(w, \phi(u, v))\) （在 \(G_{\alpha + \beta + \gamma}\) 中）。于是关系式 (19) 是下列恒等式 (《代数》， 第 I 章， § 6， no. 2， 公式 (15)) 的推论：

\[
\left(x ^ {y}, (y, z)\right). \left(y ^ {z}, (z, x)\right). \left(z ^ {x}, (x, y)\right) = e.\tag{20}
\]

在命题 1 中定义的李代数 \(\mathrm{gr}(\mathbf{G})\) （环 \(\mathbf{Z}\) 上的）称为与滤过群 \(\mathbf{G}\) 相伴的分次李代数。

